\chapter{방법}
\label{ch:method}

\section{전체 파이프라인}
\label{sec:pipeline}

한 환자에 대한 처리는 전처리와 재생으로 나뉜다. 전처리는 다음 순서로 진행된다.

\begin{enumerate}[leftmargin=2.2em]
  \item 첫 프레임에서 제공된 ED label 또는 내부 seed로 3D Chan--Vese를 풀어
  $\lev_0$을 얻는다 (제 \ref{sec:spacing-aware}절).
  \item $\surf_0$ 위에 2D Gaussian 원반을 배치하고 MRI 밝기, 불투명도, scale을
  맞추어 정준 2DGS $\Gcan$을 만든다 (제 \ref{sec:canonical}절).
  \item 프레임 $t \ge 1$에서 $\lev_{t-1}$을 seed로 Chan--Vese를 짧게 수행하여
  $\levt$를 구한다 (제 \ref{sec:warmstart}절).
  \item 각 surfel 중심을 $\surft$의 법선 방향으로 정사영하고 (제
  \ref{sec:projection}절), 두 접선축을 새 접평면으로 전송한다 (제
  \ref{sec:transport}절).
  \item 표면 이동만으로 설명되지 않는 밝기 차이를 작은 잔차로 fitting한다
  (제 \ref{sec:residual}절).
  \item 마지막 프레임까지 반복하고 $\store$만 저장한다 (제
  \ref{sec:storage-breakeven}절).
\end{enumerate}

재생에서는 Chan--Vese나 비선형 최적화를 다시 수행하지 않고, 저장된 표면을 이용한
중심/접평면 갱신과 원근 정합 래스터화만 수행한다 (제 \ref{sec:raster}절).

\paragraph{계산 비용의 구조.}
각 프레임에서 전체 2DGS를 처음부터 최적화하는 비용을 $C_{\mathrm{full}}$이라 하면
독립 최적화의 총비용은
\begin{equation}
  C_{\mathrm{independent}} = T\,C_{\mathrm{full}}
  \label{eq:cost-independent}
\end{equation}
이다. 제안 방법은 첫 프레임에서만 전체 fitting을 수행하므로
\begin{equation}
  C_{\mathrm{ours}} = C_{\mathrm{full}}
  + \sum_{t=1}^{T-1}\Bigl(C^t_{\mathrm{CV}} + C^t_{\mathrm{project}}
  + C^t_{\mathrm{orient}} + C^t_{\mathrm{res}}\Bigr)
  \label{eq:cost-ours}
\end{equation}
이다. 본 방법이 실제로 빠르려면
\begin{equation}
  C^t_{\mathrm{CV}} + C^t_{\mathrm{project}} + C^t_{\mathrm{orient}}
  + C^t_{\mathrm{res}} \;<\; C_{\mathrm{full}}
  \label{eq:cost-hypothesis}
\end{equation}
이어야 한다.

\caveat{%
식 \eqref{eq:cost-hypothesis}은 \textbf{검증할 가설이며 미리 참이라고 가정하는
결론이 아니다}. 네 항 각각을 별도로 측정하여 항별로 확인한다. 구현은 이 네 항을
독립된 타이머 스테이지로 계측한다.}

\section{간격 인지 3D Chan--Vese}
\label{sec:spacing-aware}

\subsection{에너지 범함수}

영상 분할의 변분적 토대는 Mumford--Shah 범함수~\citep{mumford1989_optimal}이다.
근사 함수 $u$가 경계 $C=\surft$로 나뉜 두 영역에서 각각 상수라고 제한하는
조각마다 상수 축약을 적용하면 매끄러움 항이 소멸하고 자료 충실도 항과 경계 측도
항만 남는다. 지시 함수를 $\Heps(\levt)$와 $1-\Heps(\levt)$로, 경계 측도를
$\deps(\levt)\norm{\gradh\levt}$의 적분으로 대체하면 다음을 얻는다.

\begin{definition}[Chan--Vese 에너지, 간격 인지 형태]
\label{def:cv-energy}
가중치 $\mu, \lin, \lout > 0$와 프레임 $t$의 level-set 함수 $\levt$, 영역 평균
$\cin, \cout$에 대하여
\begin{align}
  \ECV(\levt, \cin, \cout)
  &= \mu \int_\dom \deps(\levt)\,\norm{\gradh\levt}_2 \,dx \notag\\
  &\quad + \lin \int_\dom \abs{\imgt - \cin}^2\,\Heps(\levt)\,dx \notag\\
  &\quad + \lout \int_\dom \abs{\imgt - \cout}^2\,\bigl(1 - \Heps(\levt)\bigr)\,dx.
  \label{eq:cv-energy}
\end{align}
\end{definition}

첫째 항은 길이/면적 정규화이다. $\deps(\levt)\norm{\gradh\levt}$는 영도집합
$\surft$ 위의 근사 면적 측도이므로, 이 항은 표면 면적을 벌점으로 부과하여 매끄러운
경계를 유도한다. 둘째와 셋째 항은 각각 내부와 외부의 자료 충실도이다.

분할 문제는 $\min_{\levt,\cin,\cout}\ECV$로 정식화되며, 세 변수에 대한 결합 최소화를
교대 최소화로 접근한다.

\subsection{최적 영역 평균}

\begin{proposition}[최적 영역 평균]
\label{prop:region-means}
$\levt$를 고정하면 $\ECV$를 최소화하는 영역 평균은
\begin{align}
  \cin  &= \frac{\displaystyle\int_\dom \imgt\,\Heps(\levt)\,dx}
                {\displaystyle\int_\dom \Heps(\levt)\,dx},
  \label{eq:cin}\\[0.4em]
  \cout &= \frac{\displaystyle\int_\dom \imgt\,\bigl(1-\Heps(\levt)\bigr)\,dx}
                {\displaystyle\int_\dom \bigl(1-\Heps(\levt)\bigr)\,dx}
  \label{eq:cout}
\end{align}
이다.
\end{proposition}

\begin{proof}
$\cin$은 식 \eqref{eq:cv-energy}의 둘째 항에만 나타나므로
\[
  \frac{\partial \ECV}{\partial \cin}
  = \lin \int_\dom -2\,(\imgt - \cin)\,\Heps(\levt)\,dx = 0.
\]
$\lin>0$이므로 $\int \imgt \Heps(\levt)\,dx = \cin \int \Heps(\levt)\,dx$이고,
명제 \ref{prop:heaviside}(i)에 의해 $\Heps>0$이므로 $\int \Heps(\levt)\,dx>0$이다.
따라서 나누어 식 \eqref{eq:cin}을 얻는다. 둘째 미분은
$\partial^2 \ECV/\partial(\cin)^2 = 2\lin\int \Heps(\levt)\,dx > 0$이므로 최소점이다.
$\cout$도 같은 논법으로 얻는다.
\end{proof}

식 \eqref{eq:cin}, \eqref{eq:cout}은 $\levt$가 고정된 상태의 닫힌 형태 해이므로 교대
최소화의 각 단계에서 즉시 계산된다. $\Heps$가 급격한 이진 지시 함수가 아니므로
경계 근방의 복셀도 부분적으로 두 평균에 기여한다.

\subsection{Euler--Lagrange 방정식과 경사 흐름}

\begin{theorem}[Chan--Vese Euler--Lagrange 방정식과 경사 흐름]
\label{thm:gradient-flow}
$\cin, \cout$을 고정하고 $\levt \in C^2(\dom)$에서 $\ECV$의 최소화를 고려하자.
$\ECV$의 정류점은 다음 Euler--Lagrange 조건을 만족한다.
\begin{equation}
  \deps(\levt)\left[
    \mu\,\divh\!\left(\frac{\gradh\levt}{\norm{\gradh\levt}}\right)
    - \lin (\imgt - \cin)^2 + \lout (\imgt - \cout)^2
  \right] = 0
  \quad \text{in } \dom,
  \label{eq:euler-lagrange}
\end{equation}
자연/Neumann 경계 조건
\begin{equation}
  \frac{\deps(\levt)}{\norm{\gradh\levt}}\,\frac{\partial \levt}{\partial \nrm} = 0
  \quad \text{on } \partial\dom
  \label{eq:neumann}
\end{equation}
와 함께 성립한다. 이에 대응하며, 경사가 작은 곳의 수치 안정성을 위해 분모를
정규화한 경사 하강 흐름은 인위적 시간 $s$에 대하여
\begin{equation}
  \frac{\partial \levt}{\partial s}
  = \deps(\levt)\left[
    \mu\,\divh\!\left(\frac{\gradh\levt}{\norm{\gradh\levt}+\varepsilon}\right)
    - \lin (\imgt - \cin)^2 + \lout (\imgt - \cout)^2
  \right]
  \label{eq:gradient-flow}
\end{equation}
이다.
\end{theorem}

\begin{proof}
$\levt$의 임의의 매끄러운 섭동 $\psi \in C^\infty_c(\dom)$ 방향으로의 일차 변분을
계산한다. $E(\tau) := \ECV(\levt + \tau\psi, \cin, \cout)$로 두고
$\frac{d}{d\tau}E(\tau)\big|_{\tau=0}$을 세 항별로 계산한다. 이하 $\lev := \levt$로
줄여 쓰고, 연속 수준의 연산자를 $\nabla,\operatorname{div}$로 표기한다.

\textbf{(1) 길이 항.} $L(\lev) = \mu\int_\dom \deps(\lev)\norm{\nabla\lev}\,dx$에 대하여
$\frac{d}{d\tau}\deps(\lev+\tau\psi)\big|_0 = \deps'(\lev)\psi$이고
$\frac{d}{d\tau}\norm{\nabla(\lev+\tau\psi)}\big|_0 = \frac{\nabla\lev\cdot\nabla\psi}{\norm{\nabla\lev}}$
이므로
\begin{equation}
  \frac{dL}{d\tau}\bigg|_0
  = \mu \int_\dom \left[
    \deps'(\lev)\,\psi\,\norm{\nabla\lev}
    + \deps(\lev)\frac{\nabla\lev\cdot\nabla\psi}{\norm{\nabla\lev}}
  \right] dx.
  \label{eq:proof-length}
\end{equation}
둘째 피적분항에 벡터장 $V = \deps(\lev)\frac{\nabla\lev}{\norm{\nabla\lev}}$에 대한
발산 정리를 적용한다. $\operatorname{div}(\psi V) = \nabla\psi\cdot V + \psi\operatorname{div}V$이므로
\[
  \int_\dom \nabla\psi \cdot V\,dx
  = \int_{\partial\dom} \psi\,V\cdot \nrm\,dS - \int_\dom \psi\,\operatorname{div}V\,dx.
\]
곱셈 규칙과 $\nabla\lev\cdot\nabla\lev = \norm{\nabla\lev}^2$을 사용하면
\[
  \operatorname{div}V
  = \deps'(\lev)\,\norm{\nabla\lev}
  + \deps(\lev)\operatorname{div}\!\left(\frac{\nabla\lev}{\norm{\nabla\lev}}\right).
\]
이를 식 \eqref{eq:proof-length}에 대입하면 첫째 항의
$\deps'(\lev)\psi\norm{\nabla\lev}$와 $\operatorname{div}V$의 동일한 성분이 \emph{정확히
상쇄}되어
\begin{equation}
  \frac{dL}{d\tau}\bigg|_0
  = -\mu\int_\dom \psi\,\deps(\lev)\operatorname{div}\!\left(\frac{\nabla\lev}{\norm{\nabla\lev}}\right) dx
  + \mu\int_{\partial\dom} \psi\,\frac{\deps(\lev)}{\norm{\nabla\lev}}
    \frac{\partial\lev}{\partial\nrm}\,dS
  \label{eq:proof-length-final}
\end{equation}
를 얻는다.

\textbf{(2) 내부 충실도 항.}
$F_{\mathrm{in}}(\lev) = \lin\int_\dom \abs{\imgt-\cin}^2 \Heps(\lev)\,dx$에서 $\cin$은
고정이므로 $\Heps$만 미분되고, $\Heps' = \deps$이므로
\begin{equation}
  \frac{dF_{\mathrm{in}}}{d\tau}\bigg|_0
  = \lin \int_\dom \abs{\imgt-\cin}^2 \deps(\lev)\,\psi\,dx.
\end{equation}

\textbf{(3) 외부 충실도 항.}
$\frac{d}{d\tau}\bigl(1-\Heps(\lev+\tau\psi)\bigr)\big|_0 = -\deps(\lev)\psi$이므로
\begin{equation}
  \frac{dF_{\mathrm{out}}}{d\tau}\bigg|_0
  = -\lout \int_\dom \abs{\imgt-\cout}^2 \deps(\lev)\,\psi\,dx.
\end{equation}

\textbf{(4) 결합.} 세 항을 더하면 전체 변분은
\[
  \frac{dE}{d\tau}\bigg|_0
  = \int_\dom \psi\,\deps(\lev)\left[
    -\mu\operatorname{div}\!\left(\frac{\nabla\lev}{\norm{\nabla\lev}}\right)
    + \lin(\imgt-\cin)^2 - \lout(\imgt-\cout)^2
  \right] dx + (\text{경계항})
\]
이다. $\lev$가 정류점이면 모든 시험함수 $\psi$에 대해 이 값이 영이어야 하므로,
변분법의 기본 보조정리에 의해 대괄호에 $\deps(\lev)$를 곱한 것이 $\dom$ 내부에서
영이어야 한다. 부호를 정리하면 식 \eqref{eq:euler-lagrange}이다. 경계항
\eqref{eq:proof-length-final}이 모든 $\psi$에 대해 소멸하려면 식 \eqref{eq:neumann}이
성립해야 한다.

\textbf{(5) 경사 흐름.} $L^2(\dom)$ 내적에 대한 $\ECV$의 경사는 위 변분의 피적분
계수이다. 경사 하강은 $\partial_s\lev = -\delta \ECV/\delta\lev$로 정의되므로 부호를
정리하여 식 \eqref{eq:gradient-flow}을 얻는다. 연속 연산자를 간격 인지 연산자로
대체하고, 경사가 작은 곳에서 정규화된 곡률 항의 분모가 영이 되어 폭주하는 것을 막기
위해 $\norm{\gradh\lev}$를 $\norm{\gradh\lev}+\varepsilon$으로 정규화하였다.
\end{proof}

\begin{remark}[부호의 해석]
식 \eqref{eq:sign}의 규약에서 어떤 복셀이 내부($\lev>0$)에 있으면서 그 밝기가
$\cin$보다 $\cout$에 더 가까우면 대괄호가 음수가 되어 $\lev$가 감소하고, 그 복셀은
내부에서 빠져나간다. 이것이 의도한 방향이다. 부호를 하나만 뒤집으면 표면이
발산한다.
\end{remark}

\begin{remark}[$\divh,\gradh$로의 대체의 지위]
식 \eqref{eq:gradient-flow}에서 연속 연산자를 간격 인지 연산자로 대체한 것은 비등방
계량에서 성립하는 정확한 연속 항등식이 아니라, 연속 $L^2$ 경사 흐름의 \emph{간격
인지 이산화} 대응물로 이해해야 한다. 이 이산화를 다음 절에서 구체적으로 구성한다.
\end{remark}

\subsection{간격 인지 이산화: 첫 번째 기여}
\label{sec:discretization}

자기공명영상은 복셀 간격이 축마다 다른 비등방 격자에서 표본화되며, 슬라이스 간격
$h_z$가 평면 내 화소 간격 $h_x \approx h_y$보다 크다. 표준 유한차분이 이 간격을
무시하면 경사와 발산이 축 방향으로 왜곡되어 곡률 기반 정규화 항이 편향된다.

격자점 $(i,j,k)$에서의 이산 level-set 값을 $\lev_{i,j,k}$라 하고, 전진 차분을
\begin{equation}
  D^+_x \lev_{i,j,k} = \frac{\lev_{i+1,j,k}-\lev_{i,j,k}}{h_x},\quad
  D^+_y \lev_{i,j,k} = \frac{\lev_{i,j+1,k}-\lev_{i,j,k}}{h_y},\quad
  D^+_z \lev_{i,j,k} = \frac{\lev_{i,j,k+1}-\lev_{i,j,k}}{h_z},
  \label{eq:forward-diff}
\end{equation}
후진 차분을
\begin{equation}
  D^-_x \lev_{i,j,k} = \frac{\lev_{i,j,k}-\lev_{i-1,j,k}}{h_x},\quad
  D^-_y \lev_{i,j,k} = \frac{\lev_{i,j,k}-\lev_{i,j-1,k}}{h_y},\quad
  D^-_z \lev_{i,j,k} = \frac{\lev_{i,j,k}-\lev_{i,j,k-1}}{h_z}
  \label{eq:backward-diff}
\end{equation}
로 정의한다. 각 차분이 $1/h_x, 1/h_y, 1/h_z$로 조정된 점이 핵심이다. 간격 인지
경사는 중심 차분을 사용하여
\begin{equation}
  \gradh\lev_{i,j,k} =
  \left(
    \frac{\lev_{i+1,j,k}-\lev_{i-1,j,k}}{2h_x},\
    \frac{\lev_{i,j+1,k}-\lev_{i,j-1,k}}{2h_y},\
    \frac{\lev_{i,j,k+1}-\lev_{i,j,k-1}}{2h_z}
  \right),
  \label{eq:central-grad}
\end{equation}
간격 인지 발산은 벡터장 $V=(V_x,V_y,V_z)$에 대해 후진 차분을 사용하여
\begin{equation}
  \divh V_{i,j,k} =
  \frac{V_{x,i,j,k}-V_{x,i-1,j,k}}{h_x}
  + \frac{V_{y,i,j,k}-V_{y,i,j-1,k}}{h_y}
  + \frac{V_{z,i,j,k}-V_{z,i,j,k-1}}{h_z}
  \label{eq:backward-div}
\end{equation}
로 정의한다. 곡률 항은 $\kappa_h = \divh\bigl(\gradh\lev/\norm{\gradh\lev}\bigr)$로
이산화된다.

\paragraph{차분 조합의 근거.}
곡률 항에서 정규화된 경사장은 \emph{전진} 차분으로 만들고 발산은 \emph{후진} 차분으로
취한다. 이 짝은 임의의 선택이 아니다. 이산 부분적분 항등식에서 경계항을 제외한
부분이 연속 경우 식 \eqref{eq:proof-length-final}의 부호 구조를 보존하도록 한다.
반면 법선(식 \eqref{eq:normal-curvature})과 정사영(제 \ref{sec:projection}절)에는
중심 차분 \eqref{eq:central-grad}을 사용한다. 그 절단 오차가 $O(\hmax^2)$이며, 이것이
명제 \ref{prop:normal-error}에서 법선 각오차를 지배하는 항이다.

\begin{remark}[비등방성의 반영]
간격을 무시한 이산화는 사실상 $h_x=h_y=h_z=1$로 두는 것에 해당한다. cine CMR에서
$h_z > h_x \approx h_y$이므로 간격을 무시하면 $z$ 방향 차분이 실제보다 과대평가되어
표면 곡률과 정규화가 슬라이스 방향으로 왜곡된다. 식 \eqref{eq:central-grad},
\eqref{eq:backward-div}의 스케일링이 이 왜곡을 제거한다.
\end{remark}

\caveat{%
이 기여가 실제로 효과가 있는지는 \emph{측정해야 한다}. 구현은 모든 연산자에
\texttt{spacing\_aware} 스위치를 두어 간격을 무시한 버전을 함께 제공하며, 제
\ref{sec:ablation-design}절의 ablation은 등방 격자와 $h_z/h_x \approx 6.4$인 격자에서
두 버전의 Dice를 비교한다. 격차가 비등방성과 함께 커지지 않는다면 이 기여는 이론이
주장하는 일을 하지 않는 것이다.}

\subsection{이산 갱신식}

간격 인지 연산자를 식 \eqref{eq:gradient-flow}에 대입하고 인위적 시간을 $\Delta s$의
명시적 Euler 전진으로 이산화하면 반복 지수 $n$에 대한 갱신식
\begin{equation}
  \lev^{n+1}_{i,j,k} = \lev^{n}_{i,j,k}
  + \Delta s\,\deps\bigl(\lev^n_{i,j,k}\bigr)\left[
    \mu\,\divh\!\left(\frac{\gradh\lev^n}{\norm{\gradh\lev^n}+\varepsilon}\right)_{i,j,k}
    - \lin(\imgt{}_{,i,j,k}-\cin)^2 + \lout(\imgt{}_{,i,j,k}-\cout)^2
  \right]
  \label{eq:euler-update}
\end{equation}
을 얻는다. 각 반복에서 식 \eqref{eq:cin}, \eqref{eq:cout}으로 영역 평균을 갱신하고
식 \eqref{eq:euler-update}로 $\levt$를 갱신하는 교대 절차를 수렴할 때까지 반복한다.

\paragraph{보폭의 선택.}
식 \eqref{eq:euler-update}은 명시적 기법이므로 보폭을 고정하지 않고
\begin{equation}
  \Delta s = \mathrm{cfl}\cdot \frac{\min(\spacing)}{\max_{i,j,k}\abs{\partial_s \lev}}
  \label{eq:cfl}
\end{equation}
로 적응적으로 선택한다. 이는 한 반복에서 계면이 $\mathrm{cfl}\cdot\min(\spacing)$보다
많이 움직이지 않게 한다. 고정 보폭은 보조정리 \ref{lem:linear-conv}의 조건
$\alpha \le 1/L$을 조용히 위반할 수 있어 측정된 반복 횟수를 무의미하게 만든다.

\section{웜스타트: 두 번째 기여}
\label{sec:warmstart}

\subsection{설정}

각 프레임 $t$에 대하여 에너지 $E_t(\lev) := \ECV(\lev,\cin,\cout)$을 최소화하는 정상
상태 해를 $\lev^*_t$라 하자. 웜스타트 초기화란 프레임 $t$의 경사 하강을 시작할 때
초기값으로 이전 프레임의 수렴 해를 사용하는 것이다. 본 논문의 웜스타트는 이전 해
전체가 아니라 \emph{좁은 띠}로 제한된 초기값
\begin{equation}
  \lev^{(0)}_t = \lev^*_{t-1}
  \quad \text{on } \bigl\{\, x \in \dom : \abs{\lev^*_{t-1}(x)} < b \,\bigr\}
  \label{eq:warmstart-init}
\end{equation}
을 사용한다. $b>0$은 띠의 폭이며, 띠 밖에서는 부호만 보존하는 상수 외삽을 둔다.
이와 대비되는 콜드스타트는 프레임과 무관한 고정 초기값 $\lev_{\mathrm{cold}}$를
사용한다.

\begin{assumption}[작은 변화 가정]
\label{ass:small-change}
인접 프레임의 수렴 해가 서로 가까워, 어떤 작은 $\eta>0$에 대하여
\begin{equation}
  \norm{\lev^*_{t-1} - \lev^*_t}_{L^2(\dom)} \le \eta
  \label{eq:small-change}
\end{equation}
가 성립한다. 또한 $\lev^*_{t-1}$이 $\lev^*_t$ 주위의 국소 수렴 영역 안에 놓이며,
좁은 띠 제한 \eqref{eq:warmstart-init}이 이 근접성을 훼손하지 않는다고 가정한다.
\end{assumption}

식 \eqref{eq:small-change}의 좌변은 웜스타트의 초기 오차
$e^{\mathrm{warm}}_0 \le \eta$이다. 콜드스타트의 초기 오차
$e^{\mathrm{cold}}_0 = \norm{\lev_{\mathrm{cold}}-\lev^*_t}$은 일반적으로 이보다
훨씬 크다. 좁은 띠 초기화는 표면 근방에만 계산을 집중시켜 프레임당 비용을 줄이는
장치이며, 표면 근방에서 $\lev^*_{t-1}$과 $\lev^*_t$가 일치하므로 초기 오차의 크기를
바꾸지 않는다.

\subsection{도입하는 알려진 결과}

정리 \ref{thm:warmstart}은 새로운 수렴 정리를 처음부터 증명하는 것이 아니라, 경사
흐름과 경사 하강에 대한 알려진 결과를 웜스타트 설정에 적용한다. 어느 부분이 인용된
기지 결과이고 어느 부분이 본 논문의 특수화인지 명시한다.

\begin{lemma}[경사 흐름의 에너지 소산; 기지 결과]
\label{lem:dissipation}
$\lev(s)$가 경사 흐름 $\partial_s\lev = -\delta E_t/\delta\lev$의 매끄러운 해이면
\begin{equation}
  \frac{d}{ds}E_t(\lev(s))
  = -\norm{\frac{\delta E_t}{\delta \lev}(\lev(s))}^2 \le 0
  \label{eq:dissipation}
\end{equation}
이 성립하고, 등호는 정류점에서만 성립한다.
\end{lemma}

\begin{proof}
연쇄 법칙과 $L^2$ 내적을 사용하면
\[
  \frac{d}{ds}E_t(\lev(s))
  = \inner{\frac{\delta E_t}{\delta\lev}}{\partial_s \lev}
  = \inner{\frac{\delta E_t}{\delta\lev}}{-\frac{\delta E_t}{\delta\lev}}
  = -\norm{\frac{\delta E_t}{\delta\lev}}^2 \le 0.
\]
값이 영이 되는 것은 $\delta E_t/\delta\lev = 0$, 즉 정류점일 때뿐이다.
\end{proof}

\begin{assumption}[국소 Polyak--{\L}ojasiewicz 및 평활성]
\label{ass:pl}
조건 (i)은 Polyak~\citep{polyak1963_gradient}이 도입하고 Karimi
등~\citep{karimi2016_linear}이 경사 및 근접 경사법의 선형 수렴 조건으로 정리한
것이다.
해 $\lev^*_t$의 반지름 $r>0$ 근방
$\mathcal{B}_r = \{\lev : \norm{\lev-\lev^*_t}\le r\}$에서 다음을 가정한다.
\begin{enumerate}[label=(\roman*)]
  \item (국소 PL 조건) 어떤 상수 $m>0$에 대하여 모든 $\lev\in\mathcal{B}_r$에서
  \begin{equation}
    \tfrac{1}{2}\norm{\frac{\delta E_t}{\delta\lev}(\lev)}^2
    \ge m\bigl(E_t(\lev)-E_t(\lev^*_t)\bigr).
    \label{eq:pl}
  \end{equation}
  \item (국소 $L$-평활성) 어떤 상수 $L>0$에 대하여 $\delta E_t/\delta\lev$가
  $\mathcal{B}_r$에서 Lipschitz 상수 $L$로 Lipschitz 연속이다.
\end{enumerate}
\end{assumption}

\caveat{%
Chan--Vese 에너지 \eqref{eq:cv-energy}는 $\lev$에 대해 \textbf{전역적으로 볼록하지
않다}. 비선형 곡률 항과 매끄러운 Heaviside의 비볼록성 때문이다. 따라서 전역 강볼록성을
주장할 수 없다. 그러나 고립된 국소 최소점의 충분히 작은 근방에서는, $E_t$가 $C^2$이고
그곳의 이차 형식이 양의 정부호라는 일반적 조건 아래 가정 \ref{ass:pl}이 성립한다.
본 논문은 이 국소 조건을 가정으로 도입하며 그것이 성립하는 근방으로 논의를 제한한다.
가정 \ref{ass:small-change}은 바로 웜스타트 초기값이 이 근방에 들어옴을 보장하는
역할을 한다($\eta \le r$).}

\begin{lemma}[국소 선형 수렴; 기지 결과의 적용]
\label{lem:linear-conv}
가정 \ref{ass:pl}이 성립하고 경사 하강 반복
$\lev^{(k+1)} = \lev^{(k)} - \alpha\,\delta E_t/\delta\lev(\lev^{(k)})$를 고정 보폭
$0<\alpha\le 1/L$로 수행하며 반복열이 $\mathcal{B}_r$을 벗어나지 않는다고 하자.
그러면 에너지 잉여 $\Delta_k := E_t(\lev^{(k)})-E_t(\lev^*_t)$에 대하여
\begin{equation}
  \Delta_k \le \rho^k \Delta_0, \qquad \rho := 1-\alpha m \in (0,1)
  \label{eq:linear-conv}
\end{equation}
가 성립한다.
\end{lemma}

\begin{proof}
$L$-평활성으로부터 표준 하강 보조정리~\citep[Lemma 1.2.3]{nesterov2004_introductory}에
의해
\[
  E_t(\lev^{(k+1)}) \le E_t(\lev^{(k)})
  - \alpha\left(1-\tfrac{\alpha L}{2}\right)\norm{\frac{\delta E_t}{\delta\lev}(\lev^{(k)})}^2
  \le E_t(\lev^{(k)}) - \tfrac{\alpha}{2}\norm{\frac{\delta E_t}{\delta\lev}(\lev^{(k)})}^2,
\]
마지막 부등식은 $\alpha\le 1/L$이면 $1-\alpha L/2 \ge 1/2$임을 사용한다. 여기에 PL
조건 \eqref{eq:pl}을 대입하면
$\Delta_{k+1} \le \Delta_k - \alpha m \Delta_k = (1-\alpha m)\Delta_k$이고,
$0<\alpha m<1$이므로 귀납으로 식 \eqref{eq:linear-conv}을 얻는다.
\end{proof}

\subsection{웜스타트 수렴 정리}

\begin{theorem}[웜스타트 수렴과 반복 횟수 절감]
\label{thm:warmstart}
가정 \ref{ass:small-change}와 \ref{ass:pl}이 $\lev^*_t$의 근방 $\mathcal{B}_r$에서
성립한다고 하자. 경사 하강을 보폭 $0<\alpha\le 1/L$로 수행하고, 목표 허용 오차
$\mathrm{tol}>0$에 대해 에너지 잉여가 $\Delta_K \le \mathrm{tol}$이 되는 최소 반복
횟수를 $K(e_0)$라 하자. 여기서 $e_0 = \norm{\lev^{(0)}-\lev^*_t}$이다. 그러면 다음이
성립한다.
\begin{enumerate}[label=(\roman*)]
  \item (수렴 보존) 초기값이 $\mathcal{B}_r$ 안에 있으면 경사 하강은 $\lev^*_t$로
  선형 속도로 수렴하며, 에너지는 식 \eqref{eq:dissipation}에 의해 단조 감소한다.
  \item (정량적 반복 횟수 한계) 초기 에너지 잉여를 $L$-평활성으로 위로 어림하면
  $\Delta_0 \le \tfrac{L}{2}e_0^2$이고, 식 \eqref{eq:linear-conv}로부터
  \begin{equation}
    K(e_0) = \left\lceil \frac{\log(\Delta_0/\mathrm{tol})}{\log(1/\rho)} \right\rceil
    \le \left\lceil \frac{\log\bigl(L e_0^2/(2\,\mathrm{tol})\bigr)}{\log(1/\rho)} \right\rceil,
    \qquad \rho = 1-\alpha m.
    \label{eq:iteration-bound}
  \end{equation}
  \item (웜스타트의 이득) $e^{\mathrm{warm}}_0 \le \eta < e^{\mathrm{cold}}_0$이면
  식 \eqref{eq:iteration-bound}의 상계가 $e_0$의 단조증가 함수이므로
  \begin{equation}
    K(e^{\mathrm{warm}}_0) \le K(e^{\mathrm{cold}}_0).
    \label{eq:warmstart-gain}
  \end{equation}
\end{enumerate}
\end{theorem}

\begin{proof}
\textbf{(i)} 가정 \ref{ass:small-change}에서 $\eta \le r$이면 웜스타트 초기값
$\lev^{(0)}=\lev^*_{t-1}$은 $\mathcal{B}_r$ 안에 있다. 보조정리
\ref{lem:linear-conv}에 의해 에너지 잉여가 $\rho^k$로 감소하므로 $\Delta_k\to0$이고,
PL 조건에서 $E_t(\lev)-E_t(\lev^*_t)\ge0$이 최소점에서만 영이 되므로
$\lev^{(k)}\to\lev^*_t$이다.

\textbf{(ii)} $\delta E_t/\delta\lev(\lev^*_t)=0$이고 $\delta E_t/\delta\lev$가
$\mathcal{B}_r$에서 $L$-Lipschitz이므로 표준 이차 상계
\[
  E_t(\lev)-E_t(\lev^*_t)
  \le \inner{\frac{\delta E_t}{\delta\lev}(\lev^*_t)}{\lev-\lev^*_t}
  + \frac{L}{2}\norm{\lev-\lev^*_t}^2
  = \frac{L}{2}\norm{\lev-\lev^*_t}^2
\]
에서 $\lev=\lev^{(0)}$로 두면 $\Delta_0 \le \tfrac{L}{2}e_0^2$을 얻는다. 다음으로
$\Delta_K \le \rho^K\Delta_0 \le \mathrm{tol}$이 되기 위한 충분조건에 로그를 취하고
$\log\rho<0$임에 유의하여 정리하면
$K \ge \log(\Delta_0/\mathrm{tol})/\log(1/\rho)$이고, 이를 넘는 최소 정수가
식 \eqref{eq:iteration-bound}이다.

\textbf{(iii)} 식 \eqref{eq:iteration-bound}의 우변은 로그 안이 $e_0^2$에 비례하므로
$e_0$의 단조증가 함수이다. 따라서 결론이 따라온다.

(i)과 (ii)는 최적화 이론의 기지 결과(보조정리 \ref{lem:dissipation},
\ref{lem:linear-conv})를 인용하여 얻었고, 초기 오차를 인접 프레임 해의 근접성
\eqref{eq:small-change}과 연결하여 반복 횟수 절감 \eqref{eq:warmstart-gain}으로
특수화한 부분이 본 논문의 기여이다.
\end{proof}

\begin{remark}[선형 수렴 대 수축]
식 \eqref{eq:iteration-bound}은 에너지 잉여가 비율 $\rho$로 감소하는 선형 수렴에 대한
한계이다. 반복열 자체가 상수 $\rho'\in(0,1)$로 수축하는 경우에도 같은 로그 논법으로
$K(e_0)\le\lceil \log(e_0/\mathrm{tol})/\log(1/\rho')\rceil$을 얻으며 결론
\eqref{eq:warmstart-gain}은 동일하게 성립한다. 어느 관점에서도 더 작은 초기 오차가
더 적은 반복을 낳는다는 정성적 결론은 변하지 않는다.
\end{remark}

\caveat{%
정리 \ref{thm:warmstart}은 \emph{순서}를 예측하며 절대 반복 횟수를 예측하지 않는다.
또한 구현에서 웜스타트를 끄면 좁은 띠 crop도 함께 꺼지므로(띠가 웜스타트로부터
정의되므로), 벽시계 시간의 격차는 두 효과가 섞인다. 반복 횟수의 격차는 초기화만을
분리한다. RQ1은 두 값을 별도로 보고한다.}

\section{정준 2D Gaussian surfel}
\label{sec:canonical}

\subsection{기하}

프레임 $t$에서 surfel anchor $\anchor^t_i$가 정사영으로 $\surft$ 위에 놓인 상태라
하자. 그 점에서의 단위 법선은 간격 인지 경사로부터
\begin{equation}
  \nrm^t_i = \frac{\gradh\levt(\anchor^t_i)}{\norm{\gradh\levt(\anchor^t_i)}+\varepsilon}
  \label{eq:unit-normal}
\end{equation}
로 정의한다. 접평면 프레임 행렬은 두 접선축을 열로 하여
\begin{equation}
  \framemat^t_i = \bigl[\,e^t_{i,1}\;\; e^t_{i,2}\,\bigr] \in \Real^{3\times2},
  \qquad (\framemat^t_i)\transp \framemat^t_i = I_2,
  \qquad (\framemat^t_i)\transp \nrm^t_i = 0
  \label{eq:frame}
\end{equation}
을 만족하도록 구성한다. scale 행렬은
$\scalemat_i = \operatorname{diag}(s_{i,1},s_{i,2}) \in \Real^{2\times2}$이고, 국소
좌표 $u\in\Real^2$에 대응하는 표면 위의 점과 Gaussian 가중치는
\begin{equation}
  X^t_i(u) = \anchor^t_i + \framemat^t_i \scalemat_i u,
  \qquad
  G_i(u) = \exp\!\left(-\tfrac{1}{2}u\transp u\right)
  \label{eq:surfel-point}
\end{equation}
로 주어진다. 정의에 의해 $\{\nrm^t_i, e^t_{i,1}, e^t_{i,2}\}$은 $\anchor^t_i$에서의
오른손 정규직교 프레임을 이룬다.

식 \eqref{eq:surfel-point}에는 법선방향 두께나 3차원 공분산이 없다. 따라서
Chan--Vese의 2차원 곡면과 primitive의 차원이 직접 대응한다.

\subsection{초기 배치와 scale}

anchor는 $\surf_0$의 marching tetrahedra mesh에서 \emph{면적 균일}하게 추출한다.
면적에 비례하는 확률로 삼각형을 고르고 그 안에서 균일하게 점을 뽑으므로, 표면 위의
밀도가 고른 표본이 된다. 이후 식 \eqref{eq:projection-iter}의 정사영으로 anchor를
$\{\lev_0=0\}$ 위로 조인다. 즉 mesh는 \emph{표본 추출기}일 뿐이며 저장 표현의
일부가 아니다.

면적 $A$를 $N$개 원반이 덮을 때 평균 anchor 간격은 $\approx\sqrt{A/N}$이므로 기본
반경을
\begin{equation}
  s = \gamma\sqrt{A/N}
  \label{eq:scale-init}
\end{equation}
로 둔다. $\gamma$는 hole과 overlap 사이의 절충을 정하는 계수이며, 두 값을 모두
측정한다.

곡률 적응 확장에서는 명제 \ref{prop:planar-disk}의 평면 원반 오차
$\tfrac12\kappa s^2$이 예산 이하가 되도록 반경을 추가로 제한한다.
\begin{equation}
  s_i \le \sqrt{2\,\mathrm{budget}/\abs{\kappa_i}}.
  \label{eq:curvature-scale}
\end{equation}
이는 오차 분석을 직접적, 정량적으로 사용하는 것이다. 원반은 표면이 휘는 곳에서만
작아지고 다른 곳에서는 작아지지 않는다.

\section{법선 정사영: 세 번째 기여}
\label{sec:projection}

\subsection{최소 법선 보정의 유도}

이전 프레임의 anchor $\anchor^{t-1}_i$를 현재 표면 $\surft$로 옮기는 문제를
생각한다. 형상이 프레임 간에 조금 변하므로 이전 anchor는 일반적으로 $\surft$ 위에
없고 $\levt(\anchor^{t-1}_i)\ne0$이다.

\begin{theorem}[최소 법선 보정과 반복 투영]
\label{thm:projection}
$\levt$가 $\anchor := \anchor^{t-1}_i$ 근방에서 $C^2$이고 $\gradh\levt(\anchor)\ne0$
이라 하자. $\anchor$에서 표면 $\surft$로 향하는, 경사 방향에 평행한 일차 보정은
정규화 상수 $\varepsilon\ge0$에 대하여
\begin{equation}
  d^t_i = -\frac{\levt(\anchor^{t-1}_i)}
                {\norm{\gradh\levt(\anchor^{t-1}_i)}^2 + \varepsilon}
          \,\gradh\levt(\anchor^{t-1}_i)
  \label{eq:minimal-step}
\end{equation}
로 주어진다. 이를 반복하면 Newton 유사 투영 반복
\begin{equation}
  \anchor^{t,(k+1)}_i = \anchor^{t,(k)}_i
  - \frac{\levt\bigl(\anchor^{t,(k)}_i\bigr)}
         {\norm{\gradh\levt\bigl(\anchor^{t,(k)}_i\bigr)}^2+\varepsilon}
    \,\gradh\levt\bigl(\anchor^{t,(k)}_i\bigr),
  \qquad \anchor^{t,(0)}_i = \anchor^{t-1}_i
  \label{eq:projection-iter}
\end{equation}
을 얻으며, 이 반복은 $\surft$ 근방에서 잔차 $\levt(\anchor^{(k)})$를 영으로 국소
수렴시킨다.
\end{theorem}

\begin{proof}
$\levt$를 $\anchor$ 주위에서 일차 Taylor 전개하면 변위 $d$에 대하여
\begin{equation}
  \levt(\anchor+d) = \levt(\anchor) + \gradh\levt(\anchor)\transp d + o(\norm{d}).
  \label{eq:taylor}
\end{equation}
표면 위에 놓이는 조건 $\levt(\anchor+d)=0$을 일차항까지 부과하면
\begin{equation}
  \levt(\anchor) + \gradh\levt(\anchor)\transp d = 0.
  \label{eq:linear-constraint}
\end{equation}
식 \eqref{eq:linear-constraint}은 $d\in\Real^3$에 대한 \emph{하나의} 스칼라 방정식이므로
과소 결정된다. 표면의 형상만으로는 접선 방향 이동이 결정되지 않기 때문이다. 따라서
유일성을 위해 최소 노름 보정을 택한다.
\begin{equation}
  d = \argmin\bigl\{\,\norm{d'} \;:\; \gradh\levt(\anchor)\transp d' = -\levt(\anchor) \,\bigr\}.
  \label{eq:least-norm}
\end{equation}
임의의 허용해 $d'$을 경사 성분과 접선 성분으로
$d' = \beta\,\gradh\levt(\anchor) + w$, $w\transp\gradh\levt(\anchor)=0$로 분해하면
제약은 $\beta$만으로 결정되고 노름은
$\norm{d'}^2 = \beta^2\norm{\gradh\levt(\anchor)}^2 + \norm{w}^2$이므로 $\norm{w}$는
목적함수만 키운다. 따라서 최소해는 $w=0$, 즉 경사 방향에 평행하다.
$d=\beta\,\gradh\levt(\anchor)$를 식 \eqref{eq:linear-constraint}에 대입하면
\[
  \levt(\anchor) + \beta\norm{\gradh\levt(\anchor)}^2 = 0
  \;\Longrightarrow\;
  \beta = -\frac{\levt(\anchor)}{\norm{\gradh\levt(\anchor)}^2}.
\]
이로써 $\varepsilon=0$인 최소 법선 보정을 얻고, 경사가 작은 곳에서 분모가 영에
가까워져 보정이 폭주하는 것을 막기 위해 $\varepsilon\ge0$을 더하면 식
\eqref{eq:minimal-step}이 된다. 이 보정은 식 \eqref{eq:unit-normal}의 법선 방향에
평행하므로 법선 정사영이다.

반복 투영: 일차 보정은 일차 Taylor 근사에 기반하므로 곡률이 있는 곳에서는 한 번의
보정으로 정확히 $\surft$에 닿지 않는다. 보정을 반복하면 식
\eqref{eq:projection-iter}을 얻는다. 잔차의 이차 감소는 보조정리
\ref{lem:quadratic}에서 정량적으로 보이며, 그로부터 국소 수렴이 따라온다.
\end{proof}

\begin{proposition}[부호 거리 근방에서의 1--2회 수렴]
\label{prop:newton-steps}
$\levt$가 $\surft$ 근방에서 부호 거리 함수에 가까워 $\norm{\gradh\levt}\approx1$이
성립하면, $\varepsilon$이 무시할 만큼 작을 때 반복 \eqref{eq:projection-iter}의 한
단계는
\begin{equation}
  \anchor^{(k+1)} = \anchor^{(k)} - \levt\bigl(\anchor^{(k)}\bigr)\gradh\levt\bigl(\anchor^{(k)}\bigr)
  \label{eq:newton-step}
\end{equation}
로 축약되며, 이는 스칼라 함수 $\levt$에 대한 표준 Newton 단계와 일치한다. $\levt$가
국소적으로 거의 아핀이면 잔차는 한 번의 반복으로 거의 영이 되고, 곡률에 의한 이차
잔차는 다음 한 번의 반복에서 제거된다. 따라서 반복 횟수 $K_p$는 1 또는 2로 충분하다.
\end{proposition}

\begin{proof}
부호 거리 근방에서 $\norm{\gradh\levt}^2\approx1$이므로 $\varepsilon\to0$이면 식
\eqref{eq:projection-iter}의 분모가 1이 되어 식 \eqref{eq:newton-step}이 된다.
$\levt$가 국소적으로 아핀이면 식 \eqref{eq:taylor}의 $o(\norm{d})$ 항이 소멸하여 식
\eqref{eq:linear-constraint}이 정확한 방정식이 되고 한 번의 단계로
$\levt(\anchor^{(1)})=0$이 된다. 아핀이 아닌 경우 남는 잔차는 보조정리
\ref{lem:quadratic}에 의해 이차이므로 부호 거리 근방에서 매우 작아 추가 한 번의
반복으로 제거된다.
\end{proof}

이 성질이 구현에서 $\levt$를 주기적으로 부호 거리 함수로 재초기화하는 이유이다.
재초기화는 영도집합을 움직이지 않으면서 $\norm{\gradh\levt}\approx1$을 회복시킨다.

\subsection{접선 방향의 미결정성}
\label{sec:tangential}

식 \eqref{eq:least-norm}의 최소 노름 선택은 \emph{가장 가까운 표면으로 가는 기하학적
선택}이며 실제 심근세포 하나를 추적하는 방법이 아니다. 표면을 따라 미끄러지는 접선
운동은 두 표면의 모양만 보고는 결정할 수 없기 때문이다.

\caveat{%
그러므로 본 연구의 결과는 ``심장 모양의 시간 변화''에는 사용할 수 있으나, 추가 검증
없이 ``실제 조직의 strain''이라고 해석할 수 없다. vertex correspondence가 필요한
기능적 계측에는 mesh 운동 추적이 더 적합하다.}

또한 식 \eqref{eq:projection-iter}은 저장된 변형장이 아니라 각 프레임에서 표면이
제공하는 국소 계산이다. 따라서 위상이 바뀌거나 서로 가까운 두 표면 분기가 존재하면
잘못된 곳으로 투영될 수 있다. 구현은 최대 이동거리 제한(trust region)과 잔차가
증가한 anchor의 되돌림을 사용하여 이 실패를 줄이고, 두 사건의 발생 빈도를 보고한다.

\section{접평면 프레임 전송}
\label{sec:transport}

단위 법선 $\nrm^t_i$는 접평면을 결정하지만 그 평면 안에서의 축의 방향, 즉 면내
회전까지 고정하지는 못한다. 이 gauge 자유도를 프레임 간 연속적으로 결정하기 위해,
이전 접선축을 새 접평면으로 정사영한 뒤 정규화하는 최소 회전 전송을 사용한다.
\begin{align}
  \bar e^t_{i,1} &= e^{t-1}_{i,1} - \bigl(e^{t-1}_{i,1}\cdot \nrm^t_i\bigr)\nrm^t_i,
  \label{eq:transport-project}\\
  e^t_{i,1} &= \frac{\bar e^t_{i,1}}{\norm{\bar e^t_{i,1}}+\varepsilon},
  \qquad
  e^t_{i,2} = \nrm^t_i \times e^t_{i,1}.
  \label{eq:transport-normalize}
\end{align}
식 \eqref{eq:transport-project}은 이전 축에서 법선 성분을 제거하여 새 접평면 안에
놓이도록 하고, 식 \eqref{eq:transport-normalize}은 이를 단위 벡터로 만든 뒤 외적으로
둘째 축을 정한다. 이렇게 하면 프레임이 필요 이상으로 회전하지 않아 접선축의 시간적
연속성이 유지된다.

이는 실제 조직의 접선 이동을 추정하는 식이 아니라 rendering orientation을
연속적으로 정하는 규칙임을 명시한다. 전송의 정확성과 퇴화는 제 \ref{sec:frame-error}절에서
분석한다.

\paragraph{밀도 변화의 보정.}
법선 투영이 반복되면 수축 영역에 surfel이 몰리고 팽창 영역에 빈틈이 생길 수 있다.
이는 표면 잔차 $\Esurf$로는 드러나지 않는다. anchor가 정확히 $\surft$ 위에 있으면서도
분포가 나쁠 수 있기 때문이다. 따라서 최근접 중심거리 분포를 별도로 측정하고, 필요한
경우 접평면 안에서만 작용하는 약한 repulsion과 제한된 densification/pruning을 적용한다.
법선 성분을 제거한 뒤 이동하고 다시 정사영하므로, 이 보정이 표면 정확도를 분포
품질과 교환하지 않는다.

\section{외형 잔차}
\label{sec:residual}

표면 이동만으로 MRI 밝기 변화를 모두 설명할 수 없으므로 surfel $i$의 프레임 $t$에서의
진폭을 기준값과 잔차의 합
\begin{equation}
  \amp^t_i = \amp^0_i + \Delta \amp^t_i
  \label{eq:amplitude-split}
\end{equation}
으로 둔다. 중심, 접평면, scale을 고정하면 alpha 가중치가 주어진 2DGS 렌더링은 진폭에
대해 선형이므로(명제 \ref{prop:linearity}), 관측 pixel vector를 $y_t$, 렌더링 행렬을
$A_t$라 할 때
\begin{equation}
  \Delta \amp^*_t = \argmin_{\Delta \amp_t}
  \norm{W\bigl(A_t(\amp^0+\Delta \amp_t) - y_t\bigr)}^2_2
  + \lambda_a \norm{\Delta \amp_t}^2_2
  + \lambda_T \norm{\Delta \amp_t - \Delta \amp_{t-1}}^2_2
  \label{eq:residual-ls}
\end{equation}
를 푼다. 첫 항은 MRI appearance를 맞추고, 둘째 항은 잔차를 작게 유지하며, 셋째 항은
프레임 간 밝기 깜빡임을 억제한다.

정칙화는 선택 사항이 아니다. 알려진 slice plane에서 보이지 않는 surfel이 많으므로
$A_t$는 큰 영공간을 가지며, 자료 항만으로는 그들의 잔차가 결정되지 않는다. Tikhonov
항이 그것을 0으로 고정한다.

저장량을 더 줄이는 확장에서는
\begin{equation}
  \Delta \amp^t_i = \sum_{\ell=1}^{r} U_{i\ell} z_{\ell t},
  \qquad r \ll \min(N,T)
  \label{eq:lowrank}
\end{equation}
와 같은 저계수 시간 잔차를 사용한다. 그 절단 오차는 명제 \ref{prop:lowrank}에서
정확히 주어진다.

\section{원근 정합 래스터화}
\label{sec:raster}

단순히 3차원 Gaussian을 화면에 아핀 투영하지 않고, 각 pixel 광선과 surfel 평면의
교점을 직접 구한다. 카메라 중심을 $\camc$, pixel $q$ 방향의 \emph{단위} 벡터를
$\raydir$라 하면 광선은
\begin{equation}
  r_q(\depthq) = \camc + \depthq\,\raydir, \qquad \depthq \ge 0
  \label{eq:ray}
\end{equation}
이다. surfel $i$는 $\anchor^t_i$를 지나고 법선 $\nrm^t_i$를 갖는 평면 위에 놓이므로,
광선이 이 평면과 만나는 깊이와 교점은
\begin{equation}
  \depthq^t_i(q) = \frac{(\nrm^t_i)\transp\bigl(\anchor^t_i - \camc\bigr)}
                        {(\nrm^t_i)\transp \raydir},
  \qquad
  x^t_i(q) = \camc + \depthq^t_i(q)\,\raydir
  \label{eq:ray-plane}
\end{equation}
로 계산한다. $\raydir$가 단위이므로 $\depthq$는 mm 단위 거리이며 복셀 간격과 직접
비교 가능하다. 교점의 2차원 국소 좌표는
\begin{equation}
  u^t_i(q) = \scalemat_i^{-1}(\framemat^t_i)\transp\bigl(x^t_i(q) - \anchor^t_i\bigr),
  \label{eq:local-coord}
\end{equation}
화면 pixel에서의 불투명도는
\begin{equation}
  \alphaq^t_i(q) = \opa^t_i \exp\!\left(-\tfrac{1}{2}\norm{u^t_i(q)}^2_2\right)
  \label{eq:alpha}
\end{equation}
이다. $(\nrm^t_i)\transp\raydir$가 0에 가까운 스침각 surfel과 카메라 뒤의 교점
($\depthq \le 0$)은 제거한다. 남은 surfel을 $\depthq^t_i(q)$ 순서로 alpha
composition하면
\begin{equation}
  \Chat_t(q) = \sum_i \left(\prod_{j<i}\bigl(1-\alphaq^t_j(q)\bigr)\right)
               \alphaq^t_i(q)\,c^t_i
  \label{eq:compositing}
\end{equation}
를 얻는다. $c^t_i$는 진폭 $\amp^t_i$를 grayscale 또는 의료영상 window/level colormap으로
변환한 값이다.

\subsection{정준 프레임 fitting의 손실}

첫 프레임의 2DGS fitting에는 세 종류의 제약을 추가한다. 투영된 Chan--Vese mask와
rendered alpha mask 사이의 $\mathcal{L}_{\mathrm{mask}}$, rendered depth에서 계산한
normal과 식 \eqref{eq:unit-normal} 사이의 $\mathcal{L}_{\mathrm{normal}}$, 같은 pixel에
기여하는 surfel들의 depth가 불필요하게 벌어지는 것을 억제하는
$\mathcal{L}_{\mathrm{dist}}$이다. 전체 손실을
\begin{equation}
  \mathcal{L} = \mathcal{L}_{\mathrm{app}}
  + \lambda_m \mathcal{L}_{\mathrm{mask}}
  + \lambda_n \mathcal{L}_{\mathrm{normal}}
  + \lambda_d \mathcal{L}_{\mathrm{dist}}
  \label{eq:total-loss}
\end{equation}
로 둔다.

\caveat{%
기하가 Chan--Vese 표면에 고정되므로 이 손실은 경계를 바꿀 수 없다. 구현은 anchor,
접선축, 법선을 최적화 불가능한 버퍼로 등록하여 이를 학습률이 아니라 \emph{구조적으로}
보장한다(제 \ref{sec:impl-decisions}절). 따라서 식 \eqref{eq:total-loss}의 기하 항들은
통상의 2DGS에서와 다른 역할을 한다. 표면을 개선할 수는 없고, 고정된 표면을 원반이
덮는 방식만 개선한다.}

\subsection{재생 시 프레임 예산}

재생 중 프레임당 계산은 표면 loading, 식 \eqref{eq:projection-iter}의 소수 갱신,
식 \eqref{eq:unit-normal}--\eqref{eq:transport-normalize}의 orientation 갱신, 식
\eqref{eq:compositing}의 래스터화이다. 목표 조건은
\begin{equation}
  T_{\mathrm{frame}} = T_{\mathrm{surface}} + T_{\mathrm{project}}
  + T_{\mathrm{orient}} + T_{\mathrm{raster}} < 33.3\,\mathrm{ms}
  \label{eq:frame-budget}
\end{equation}
이다. 네 항을 각각 계측한다.

\section{저장 표현과 손익분기}
\label{sec:storage-breakeven}

모든 프레임의 2DGS를 저장하는 경우
\begin{equation}
  \Sfull = T\,N\,P_{2D}
  \label{eq:s-full}
\end{equation}
이고, 제안 표현 \eqref{eq:store}은
\begin{equation}
  \Sours = N P_{2D} + \sum_t S(\surft) + T\,N\,P_r
  \label{eq:s-ours}
\end{equation}
이다. $P_{2D}$는 한 surfel의 parameter 수, $P_r$은 잔차 parameter 수이다. 압축비는
\begin{equation}
  \CR = \frac{\Sfull}{\Sours}
  \label{eq:cr}
\end{equation}
로 계산한다.

\paragraph{손익분기 조건.}
$\Sours = \Sfull$을 프레임당 표면 예산에 대해 풀면
\begin{equation}
  S^{\ast}(\surf) = \frac{T N P_{2D} - N P_{2D} - T N P_r}{T}
  = N P_{2D}\left(1 - \frac{1}{T}\right) - N P_r
  \label{eq:breakeven}
\end{equation}
을 얻는다. 프레임당 표면 저장량이 $S^\ast$를 넘으면 $\CR<1$이 되어 전 프레임 2DGS를
저장하는 것이 오히려 싸다. 식 \eqref{eq:breakeven}이 음수이면 잔차만으로도 이미
$\Sfull$을 초과하므로 이 구성에서는 제안 표현이 이길 수 없다.

\caveat{%
따라서 RQ3은 ``압축비가 얼마인가''가 아니라 ``애초에 1을 넘는가''를 먼저 묻는다.
구현은 표면을 mesh, binary mask, 좁은 띠 부호 거리의 세 방식으로 저장할 때의 크기를
\emph{모두} 보고하고, 이론적 바이트 수와 실제 파일 크기를 나란히 제시한다.}

\section{중간 시점 보간}
\label{sec:interpolation}

원 cine CMR이 20--40 phase라 하더라도 재생을 부드럽게 하기 위해 인접 표면의 부호
거리와 잔차를 보간할 수 있다. $\tau = t+\beta$, $0\le\beta\le1$에 대하여
\begin{equation}
  \lev_\tau = (1-\beta)\levt + \beta \lev_{t+1},
  \qquad
  \Delta \amp_\tau = (1-\beta)\Delta \amp_t + \beta \Delta \amp_{t+1}
  \label{eq:interpolation}
\end{equation}
로 두고 같은 법선 정사영을 적용한다.

\caveat{%
이는 \textbf{표시를 부드럽게 하는 보간일 뿐}이며, 관측되지 않은 임상 정보를 새로
복원하는 것으로 해석하지 않는다. 명제 \ref{prop:interp-eikonal}은 보간된 level-set이
참 중간 표면의 부호 거리 함수가 아님을 정량적으로 보이고, 명제
\ref{prop:interp-projection}은 그것이 정사영에 주입하는 오차를 명시한다. 임상적으로
유효한 중간 표면이 필요하면 그 시점의 영상을 직접 분할해야 한다.}
